\begin{longtable}{ | l | m{5cm} | l | l | }
\hline\rowcolor{lightgray}
{\bfseries 名称} & {\bfseries 描述} &  {\bfseries 地址} & {\bfseries 访问} \\ \hline
\endhead
\multicolumn{4}{|l|}{\textbf{HP CPU 总线记录配置寄存器}} \\ \hline
\hyperref[regdesc:SPMMEMMONITORLOGSETTINGREG]{SPM\_MEM\_MONITOR\_LOG\_SETTING\_REG} & 总线访问记录配置寄存器 & 0x0000 & R/W \\ \hline
\hyperref[regdesc:SPMMEMMONITORLOGCHECKDATAREG]{SPM\_MEM\_MONITOR\_LOG\_CHECK\_DATA\_REG} & 总线访问监测数据配置寄存器 & 0x0008 & R/W \\ \hline
\hyperref[regdesc:SPMMEMMONITORLOGDATAMASKREG]{SPM\_MEM\_MONITOR\_LOG\_DATA\_MASK\_REG} & 总线访问监测数据屏蔽配置寄存器 & 0x000C & R/W \\ \hline
\hyperref[regdesc:SPMMEMMONITORLOGMINREG]{SPM\_MEM\_MONITOR\_LOG\_MIN\_REG} & 总线访问监测范围配置寄存器 & 0x0010 & R/W \\ \hline
\hyperref[regdesc:SPMMEMMONITORLOGMAXREG]{SPM\_MEM\_MONITOR\_LOG\_MAX\_REG} & 总线访问监测范围配置寄存器 & 0x0014 & R/W \\ \hline
\hyperref[regdesc:SPMMEMMONITORLOGMEMSTARTREG]{SPM\_MEM\_MONITOR\_LOG\_MEM\_START\_REG} & 记录数据写入内存的起始地址 & 0x0018 & R/W \\ \hline
\hyperref[regdesc:SPMMEMMONITORLOGMEMENDREG]{SPM\_MEM\_MONITOR\_LOG\_MEM\_END\_REG} & 记录数据写入内存的结束地址 & 0x001C & R/W \\ \hline
\hyperref[regdesc:SPMMEMMONITORLOGMEMCURRENTADDRREG]{SPM\_MEM\_MONITOR\_LOG\_MEM\_CURRENT\_ADDR\_REG} & 表示下一次写内存的写地址 & 0x0020 & RO \\ \hline
\hyperref[regdesc:SPMMEMMONITORLOGMEMADDRUPDATEREG]{SPM\_MEM\_MONITOR\_LOG\_MEM\_ADDR\_UPDATE\_REG} & 更新下一次写内存的写地址为记录数据写入内存的起始地址 & 0x0024 & WT \\ \hline
\hyperref[regdesc:SPMMEMMONITORLOGMEMFULLFLAGREG]{SPM\_MEM\_MONITOR\_LOG\_MEM\_FULL\_FLAG\_REG} & 记录溢出状态寄存器 & 0x0028 & 不定 \\ \hline
\multicolumn{4}{|l|}{\textbf{时钟控制寄存器}} \\ \hline
\hyperref[regdesc:SPMMEMMONITORCLOCKGATEREG]{SPM\_MEM\_MONITOR\_CLOCK\_GATE\_REG} & 寄存器时钟控制 & 0x002C & R/W \\ \hline
\multicolumn{4}{|l|}{\textbf{版本寄存器}} \\ \hline
\hyperref[regdesc:SPMMEMMONITORDATEREG]{SPM\_MEM\_MONITOR\_DATE\_REG} & 版本寄存器 & 0x03FC & R/W \\ \hline
\end{longtable}
